\subsection*{The Whole Numbers as a Structure}
\label{subsec:whole-numbers-as-a-structure}

\begin{remark*}[Modeling Goal]
We exhibit the whole numbers as a concrete structure with carrier
$\mathbb{W}=\{0,1,2,\ldots\}$. In the usual convention this is extensionally the
same carrier as $\mathbb{N}$ when $\mathbb{N}$ is zero-based. It is kept here as
a named presentation because elementary arithmetic often distinguishes
``natural numbers'' from ``whole numbers'' terminologically.
\end{remark*}

\begin{remark*}[Exposition]
The whole-number signature is the zero-based arithmetic fragment
\[
\mathcal{L}_{\mathbb{W}}=\{S,+,\cdot,0,1\}.
\]
It has unary successor $S$, binary operations $+,\cdot$, and constants $0,1$.
The order symbol $<$ is absent at first and is added by expansion below.
\end{remark*}

\begin{definition}[The Whole Number Model $\mathcal{W}$]
\label{def:whole-number-model}
The \textbf{whole number model} is the $\mathcal{L}_{\mathbb{W}}$-structure
\[
\mathcal{W}
=
\bigl(\mathbb{W};\,
S^{\mathcal{W}},+^{\mathcal{W}},\cdot^{\mathcal{W}},
0^{\mathcal{W}},1^{\mathcal{W}}\bigr),
\]
where
\[
\mathbb{W}=\{0,1,2,\ldots\},
\]
$S^{\mathcal{W}}$ is successor, $+^{\mathcal{W}}$ is whole-number addition,
$\cdot^{\mathcal{W}}$ is whole-number multiplication, and
$0^{\mathcal{W}},1^{\mathcal{W}}$ are the distinguished whole-number constants.
\end{definition}

\begin{remark*}[Standard quantified statement]
\[
|\mathcal{W}|=\mathbb{W},\qquad
S^{\mathcal{W}}:\mathbb{W}\to\mathbb{W},\qquad
+^{\mathcal{W}},\cdot^{\mathcal{W}}:\mathbb{W}\times\mathbb{W}\to\mathbb{W},
\qquad
0^{\mathcal{W}},1^{\mathcal{W}}\in\mathbb{W}.
\]
\end{remark*}

\begin{remark*}[Interpretation]
If the repository defines $\mathbb{N}$ as $\{0,1,2,\ldots\}$, then
$\mathcal{W}$ is not a new mathematical structure; it is the same zero-based
counting structure under the name ``whole numbers.'' If the repository defines
$\mathbb{N}$ as $\{1,2,3,\ldots\}$, then $\mathbb{W}$ is the zero-adjoined
successor structure.
\end{remark*}

\begin{remark*}[Satisfaction]
The whole-number model satisfies the same zero-based Peano and arithmetic
laws as the natural-number model whenever the repository convention identifies
\[
\mathbb{W}=\mathbb{N}.
\]
Equivalently,
\[
\mathcal{W}\models \Sigma_{\mathbb{W}},
\qquad
\operatorname{PeanoSystem}(\mathbb{W},S,0).
\]
This satisfaction claim should be wired to the repository theorem that verifies
the Peano-system or commutative-semiring laws for the zero-based counting
structure.
\end{remark*}

\begin{remark*}[Expansion to the ordered whole-number model]
Adjoining the relation symbol $<$ gives the expansion
\[
\mathcal{W}_{<}=(\mathbb{W};S,+,\cdot,0,1,<),
\qquad
\operatorname{Expansion}(\mathcal{W}_{<},\mathcal{W}),
\]
where $<^{\mathcal{W}}$ is the usual order on $\{0,1,2,\ldots\}$.
\end{remark*}

\begin{remark*}[Reduct to the successor structure]
Forgetting $+,\cdot,1$ leaves the successor reduct
\[
\mathcal{W}\!\restriction\!\{S,0\}
=
(\mathbb{W};S,0),
\qquad
\operatorname{Reduction}\!\bigl(\mathcal{W}\!\restriction\!\{S,0\},\,\mathcal{W}\bigr).
\]
This is the zero-based counting structure.
\end{remark*}

\begin{remark*}[Reduct to the additive monoid]
Forgetting $S,\cdot,1$ leaves
\[
\mathcal{W}\!\restriction\!\{+,0\}
=
(\mathbb{W};+,0),
\]
the additive commutative monoid of whole numbers.
\end{remark*}

\begin{remark*}[Blueprint position]
The whole numbers occupy a terminological rather than algebraically new
position. They emphasize that zero is included:
\[
\mathbb{W}=\{0,1,2,\ldots\}.
\]
Structurally, they sit at the zero-based counting layer and embed into the
integers:
\[
\mathbb{W}\hookrightarrow\mathbb{Z}.
\]
Passing from $\mathbb{W}$ to $\mathbb{Z}$ supplies additive inverses.
\end{remark*}

\begin{remark*}[Interpretation]
The whole numbers are closed under addition and multiplication, contain $0$
and $1$, and support successor. They are not a group, ring, or field under the
usual arithmetic operations because additive inverses and multiplicative inverses
are not available internally. Their natural algebraic home is the same
commutative-semiring / zero-based Peano layer as the natural-number model.
\end{remark*}

